

\chapter{The ponderomotive force theory}
\section{Equations of motion}

Consider the relativistic equations of motion for a particle with a time and coordinate dependent mass ${\cal M}({\bm x},t)$

The covariant Lagrange function will be
\begin{align}
 L_{(\tau)} = -{\cal M}(x) c\sqrt{\frac{dx_\mu}{d\tau}\frac{dx^\mu}{d\tau}}=-{\cal M}(x) c\sqrt{\dot x_{\mu}\dot x^{\mu}}
\end{align}
The equations of motion are obtained from
\begin{align}
 &\frac{\partial L_{(\tau)}}{\partial\dot x_\mu}=-{\cal M}(x)c\frac{\dot x^{\mu}}{\sqrt{\dot x_{\mu}\dot x^{\mu}}}\quad \frac{\partial L_{(\tau)}}{\partial x_{\mu}}=-\frac{\partial {\cal M}(x)}{\partial x_\mu}c\sqrt{\dot x_{\mu}\dot x^{\mu}},\\
  &\frac{d}{d\tau}\left(-{\cal M}(x)c\frac{\dot x^{\mu}}{\sqrt{\dot x_{\mu}\dot x^{\mu}}}\right)=-\frac{\partial {\cal M}(x)}{\partial x_{\mu}}c\sqrt{\dot x_\mu\dot x^\mu}
\end{align}
If we assume that
\begin{align}
 \dot x_\mu\dot x^{\mu} = c^2
\end{align}
as it should be then the equations of motion are
\begin{align}
    \frac{d}{d\tau}\left({\cal M}(x)\dot x^{\mu}\right)=c^2\frac{\partial {\cal M}(x)}{\partial x_{\mu}}
\end{align}
The expression of the four acceleration $\frac{du}{d\tau}$ will be  
\begin{align}
{\cal M}(x)\frac{d\dot x^{\mu}}{d\tau}=c^2\frac{\partial{\cal M}(x)}{\partial x_{\mu}}-\dot x^{\mu} \frac{d{\cal M}(x)}{d\tau}=c^2\frac{\partial{\cal M}(x)}{\partial x_{\mu}}-\dot x^{\mu}\dot x^{\nu}\frac{\partial{\cal M}(x)}{\partial x_{\nu}}=\left(c^2 g^{\mu\nu} -\dot x^{\mu}\dot x^{\nu}\right)\frac{\partial{\cal M}(x)}{\partial x^{\nu}}
\end{align}
The covariant form of the equations of motion 
\begin{align}
    \boxed{\frac{dx^{\mu}}{d\tau}=u^{\mu},\quad \frac{du^{\mu}}{d\tau}=\frac{1}{{\cal M}(x)}\left(c^2 g^{\mu\nu} -\dot x^{\mu}\dot x^{\nu}\right)\frac{\partial{\cal M}(x)}{\partial x^{\nu}}}
\end{align}
To check the correctness of the above formula we shoud verify that
\begin{align}
    u\cdot \frac{du}{d\tau}=0
\end{align}
In our case,
\begin{align}
u_\mu\frac{du^\mu}{d\tau}=\frac{1}{{\cal M}(x)}\left(c^2 g^{\mu\nu} -u^{\mu}u^{\nu}\right)u_\mu\frac{\partial{\cal M}(x)}{\partial x^{\nu}}=0
\end{align}


The above formulas are valid if $u^\mu\approx 0$, or $u^\mu\ll c $. If $u$ is not small, the formulas are still valid, but in a reference frame in which the dressed electron is at rest, and we need to perform a Lorentz boost to find the acceleration in the laboratory frame.





\section{The dressed mass}

In our case we define the variable mass in terms of a 4-vector potential
\begin{align}
 A(x) = \left(\frac{ \Phi(x)}c,{\bm A}(x)\right)
\end{align}
as
\begin{align}
 {\cal M}(x) = m\sqrt{1-\frac{e^2\langle A^2\rangle (x)}{(mc)^2}}
\end{align}
where $m$ is the (constant) mass of the isolated particle and $ \langle A^2\rangle(x) $ is the time average
\begin{align}
 \langle A^2\rangle (x)= \frac 1{T} \int\limits_{-T/2}^{T/2}d\tau A^2(t+\tau, {\bm x})
\end{align}
In  a Coulomb gauge, the four vector potential reduces to
\begin{align}
 A(x) = \left(0, {\bm A}(x)\right), \quad A^2(x) = -{\bm A}^2(x)
\end{align}
and the dressed mass becomes
\begin{align}
 {\cal M}(x) =m\sqrt{1+\frac{e^2\langle {\bm A}^2\rangle(x)}{(mc)^2}},\quad \langle {\bm A}^2\rangle (x)= \frac 1{T} \int\limits_{-T/2}^{T/2}d\uptau {\bm A}^2(t+\uptau, {\bm x}) = -\frac 1{T} \int\limits_{-T/2}^{T/2}d\uptau { A}^2(t+\uptau, {\bm x})
\end{align}
Define the dimensionless quantity
\begin{align}
 a(x) = \frac{e^2}{(mc)^2}\frac1T\int\limits_{-T/2}^{T/2} d \uptau {\bm A}^2(t+\uptau, {\bm x})= -\frac{e^2}{(mc)^2}\frac1T\int\limits_{-T/2}^{T/2} d \uptau { A}^2(t+\uptau, {\bm x})
\end{align}

The derivatives are
\begin{align}
 \frac{\partial a(x)}{\partial x_{\mu}}= \frac{2e^2}{(mc)^2}\frac1T\int\limits_{-T/2}^{T/2} d \uptau {\bm A}(t+\uptau, {\bm x})\cdot \partial^\mu {\bm A}(t+\uptau, {\bm x})=- \frac{2e^2}{(mc)^2}\frac1T\int\limits_{-T/2}^{T/2} d \uptau { A}(t+\uptau, {\bm x})\cdot \partial^\mu { A}(t+\uptau, {\bm x})
\end{align}


We can define the 4-vector
\begin{align}
 \upsilon = (1,0,0,0),\quad  \text{and}\quad c\uptau\upsilon = (c\uptau, 0, 0, 0)
\end{align}

and define ${\bm A}$ as a function of the four-vector $x\equiv (ct,{\bm x})$. Then in covariant notations we get
\begin{align}
& \boxed{a(x) = -\frac{e^2}{(mc)^2}\frac1{\lambda}\int\limits_{-\lambda/2}^{\lambda/2} d\chi  { A}^2(x + \chi \upsilon)=-\frac{e^2}{(mc)^2}\frac1{\lambda}\int\limits_{-\lambda/2}^{\lambda/2} d\chi   A_\alpha(x + \chi \upsilon) A^\alpha(x + \chi \upsilon)}\\
&\boxed{\partial^\mu a(x) = -\frac{2e^2}{(mc)^2}\frac1{\lambda}\int\limits_{-\lambda/2}^{\lambda/2} d\chi  { A}(x + \chi \upsilon)\cdot\partial^\mu{ A}(x + \chi \upsilon)=-\frac{2e^2}{(mc)^2}\frac1{\lambda}\int\limits_{-\lambda/2}^{\lambda/2} d\chi  { A}_\alpha(x + \chi \upsilon)\cdot\partial^\mu{ A}^\alpha(x + \chi \upsilon)}\\
&\boxed{{\cal M}(x) = m  \sqrt{1+a(x)}}\\
&\boxed{\partial^\mu {\cal M}(x) =m\partial^\mu\sqrt{1+a(x)}=\frac m2{\partial^\mu a(x)}\frac1{\sqrt{1+a(x)}}}
\end{align}
